\chapter{Tabla ampliada de parámetros y relevamiento}
\label{app:parametros}

Este apéndice reúne parámetros que deberían mantenerse trazables durante el proyecto. Algunos ya aparecen en la simulación actual; otros pertenecen a una etapa de laboratorio o campus. La tabla permite distinguir entre valores usados, valores a medir y valores a decidir.

\section{Parámetros físicos}

\begin{longtable}{p{0.20\textwidth}p{0.16\textwidth}p{0.22\textwidth}p{0.30\textwidth}}
  \caption{Parámetros físicos relevantes para el enlace.}
  \label{tab:parametros_fisicos_ampliados}\\
  \toprule
  Parámetro & Símbolo & Estado actual & Comentario \\
  \midrule
  \endfirsthead
  \toprule
  Parámetro & Símbolo & Estado actual & Comentario \\
  \midrule
  \endhead
  Distancia de fibra & \(L\) & Barrida en simulación & Reemplazar por longitud óptica medida en campus \\
  Atenuación & \(\alpha\) & \SI{0.2}{\decibel\per\kilo\meter} & Valor típico; medir si se accede a fibra real \\
  Pérdida de conectores & \(A_c\) & No modelada explícitamente & Puede agregarse al presupuesto de enlace \\
  Número medio de fotones & \(\mu\) & 0,1 / 0,6 según experimento & Debe calibrarse a la salida de Alice \\
  Intensidad señuelo & \(\nu\) & 0,2 en Exp. 4 & En decoy real se elige aleatoriamente por pulso \\
  Frecuencia de repetición & \(f_\text{rep}\) & \SI{80}{\mega\hertz} & Limitada por costo computacional de la simulación \\
  Visibilidad & \(V\) & Barrida en Exp. 3 & Medir en interferómetro real \\
  Error de fase & \(p_\text{fase}\) & \((1-V)/2\) & Abstracción del modelo \\
  Eficiencia de detector & \(\eta_d\) & Barrida en Exp. 2 & Depende de tecnología y operación \\
  Conteos oscuros & \(DCR\) & Barridos en Exp. 2 & Deben normalizarse por ventana temporal \\
  Resolución temporal & \(\Delta t\) & \SI{10}{\pico\second} & Impacta asignación de bins \\
  Tiempo muerto & \(t_d\) & Modelado indirecto & Relevante a tasas altas \\
  \bottomrule
\end{longtable}

\section{Parámetros de protocolo}

\begin{longtable}{p{0.22\textwidth}p{0.24\textwidth}p{0.42\textwidth}}
  \caption{Parámetros de protocolo y postprocesamiento.}
  \label{tab:parametros_protocolo}\\
  \toprule
  Parámetro & Estado actual & Próximo paso \\
  \midrule
  \endfirsthead
  \toprule
  Parámetro & Estado actual & Próximo paso \\
  \midrule
  \endhead
  Elección de bases & Implementada por BB84 de SeQUeNCe & Documentar probabilidades y sesgos si se modifican \\
  Longitud de llave & 128 bits por muestra & Aumentar para reducir varianza estadística \\
  Número de llaves & 3 a 5 por punto & Repetir con múltiples semillas \\
  Estimación de QBER & Promedio de \texttt{error\_rates} & Separar muestra de estimación y muestra de clave \\
  Corrección de errores & No implementada como protocolo & Evaluar Cascade o LDPC en etapa posterior \\
  Amplificación de privacidad & Factor de entropía en SKR & Implementar hashing universal si se genera llave final \\
  Decoy states & Cota analítica híbrida & Implementar selección aleatoria de intensidades \\
  Autenticación clásica & Fuera del modelo & Definir mecanismo en implementación real \\
  \bottomrule
\end{longtable}

\section{Checklist de relevamiento Campus Miguelete}

El checklist siguiente está pensado para una visita técnica. La salida ideal es una planilla con un registro por enlace candidato.

\begin{longtable}{p{0.08\textwidth}p{0.36\textwidth}p{0.44\textwidth}}
  \caption{Checklist de relevamiento para enlaces candidatos.}
  \label{tab:checklist_campus}\\
  \toprule
  ID & Pregunta & Evidencia esperada \\
  \midrule
  \endfirsthead
  \toprule
  ID & Pregunta & Evidencia esperada \\
  \midrule
  \endhead
  R1 & ¿Qué edificios o salas serían nodos? & Nombre, ubicación, responsable, fotos \\
  R2 & ¿Existe fibra entre esos puntos? & Plano, patch panel, disponibilidad de pelos \\
  R3 & ¿Cuál es la longitud óptica estimada? & Medición OTDR o documentación de tendido \\
  R4 & ¿Cuántos conectores/empalmes hay? & Lista de puntos de conexión \\
  R5 & ¿Qué tipo de conector se usa? & FC/PC, FC/APC, LC u otro \\
  R6 & ¿Hay rack y energía estable? & Fotos, tomas, UPS, puesta a tierra \\
  R7 & ¿La sala tiene acceso controlado? & Cerradura, registro, política de ingreso \\
  R8 & ¿Cómo es el ambiente térmico? & Temperatura, ventilación, equipos cercanos \\
  R9 & ¿Hay vibración o manipulación de fibra? & Observación de ductos, bandejas, tránsito \\
  R10 & ¿Hay conectividad IP para canal clásico? & VLAN, dirección, firewall, latencia aproximada \\
  R11 & ¿Dónde se ubicaría el KMS? & Nodo confiable, sala segura, equipo de control \\
  R12 & ¿Qué aplicación probaría las llaves? & Servidor, software, interfaz prevista \\
  \bottomrule
\end{longtable}

\section{Plantilla de ficha de enlace}

Cada enlace candidato debería documentarse con una ficha única:

\begin{itemize}
  \item \textbf{Identificación}: nombre del enlace, edificios, salas y responsables.
  \item \textbf{Ruta física}: ductos, patch panels, fibras disponibles y conectores.
  \item \textbf{Pérdidas}: longitud, atenuación estimada, pérdida fija y margen.
  \item \textbf{Condiciones de sala}: rack, energía, temperatura, acceso y seguridad.
  \item \textbf{Parámetros para simulación}: \(L\), \(\alpha\), pérdidas fijas, restricciones de detector y canal clásico.
  \item \textbf{Riesgos}: puntos de falla, permisos, mantenimiento y dependencia de terceros.
\end{itemize}

Esta ficha evita que el diseño dependa de memoria oral o supuestos. También permite comparar enlaces candidatos con una misma escala. Si dos rutas tienen distancia similar pero una tiene más conectores y peor seguridad física, la simulación y el análisis de riesgos pueden justificar elegir la otra.
